\subsection{The first two leads}

\begin{result}{A.8}{Open}{Lead 1: the four referee points}
The signed-frame identity, the four quotient comparisons,
the prime-power centralizer passage and the ordinary
pentagon dependency have been checked in A.1--A.3.
The complete independent finite census demanded by its
second point remains the unresolved task A.4.
The full-completion comparison is recorded separately
under A.6. Thus Lead 1 is not yet a completed task.
\end{result}

\begin{result}{A.9}{Proved}{Lead 2: the relevant solvable centralizer equality}
Every finite subgroup used in the E.44/E.49 chain is
a finite prime-power group. The completed centralizer
equalities follow directly from Boggi--Zalesskii
Theorem A(iii), with the additional induced-topology
proofs supplied in A.1. The solvable genus-zero
assertion is proved in A.1.2. This repairs the
citation route suggested by Lead 2: its desired
centralizer equality holds here without the
unjustified arbitrary closure-intersection step.
\end{result}
\begin{proof}
The ten-point group is $C_2^2$, the translation group
is $C_3^2$, and the remaining finite subgroups are
cyclic of order two or three. The cited prime-power
theorem applies to each complete finite subgroup,
not just to its generators. The specific hypotheses
and the full-topology maps are established in A.1
and in the labelled comparison of A.3. For a
general solvable finite subgroup in genus zero,
A.1.2 gives the induction through its characteristic
elementary-abelian subgroup. The proof there also
explains how the finite conjugation map to
$\operatorname{Aut}(H)$ passes from normalizers
to centralizers. These facts prove the asserted
equalities. They make no assertion for the
icosahedral subgroup outside this task.
\end{proof}
